\section{Experimental Setup}
\label{sec:experimental-setup}

We evaluate two provider-diverse models: \qwen{} through Together AI and \gptnew{} through \gptprovider{}. \qwenshort{} evaluation is complete; \gptnew{} evaluation is in progress and all corresponding cells are marked \tbd{}. All evaluations use temperature~0.7 with three independent samples per source-target pair (no chain-of-thought mode, single-attempt translation with no self-repair). Model and hardware details are in Appendix Tables~\ref{tab:model-config} and~\ref{tab:hardware}; token and cost accounting is in Appendix~\ref{sec:appendix-cost}.

The standard evaluation matrix covers six directions: CUDA$\leftrightarrow$OpenMP, CUDA$\leftrightarrow$OpenCL, and OpenMP$\leftrightarrow$OpenCL. Two OpenMP-target directions are included as case studies because they are available only in a smaller subset and require a different compiler stack. For \qwenshort{}, the canonical evaluation covers 142 unique L0 source-target pairs across all eight directions, with three samples each (426 L0 records). A separate augmentation campaign evaluates a subset of eligible pairs at levels L1--L4 (50 records per level, 200 total), for 626 valid records after excluding 82 records involving 9 \knownfail{} specs.

The primary metrics are pass@$k$ (Chen et al., 2021) computed over the 142 unique L0 tasks and the per-record pass rate across all 626 records. We report Wilson 95\% confidence intervals for proportions. For augmentation analysis, we report descriptive per-level pass rates on a balanced 12-kernel CUDA-to-OpenMP subset present at all five levels; the small per-level sample sizes (12 records at L1--L4) preclude formal trend tests. Performance speedup is intentionally out of scope: wall-clock measurements in the current harness conflate kernel runtime with I/O, allocation, launch overhead, and scheduling noise.
